\chapter{Category-Theoretic Interface and Terminality}

The ISAR kernel provides a category-theoretic interpretation of representation-free substrates and semantic views. We formalize this using the category of semantic kernels and morphisms, and prove that the canonical invariant quotient space is a terminal object in this category. This universal property of terminality establishes the mathematical foundation of view-independence.

\section{The Category of Semantic Kernels}

The definitions and structures in this section are mechanically verified in \url{src/ISAR/KernelCategory.lean} and \url{src/ISAR/DialectKernel.lean}. We define the category $\mathcal{K}$ of admissible semantic kernels over the substrate.

\begin{definition}[Semantic Kernel]
\lean{ISAR.Kernel}
An admissible semantic kernel $K \in \mathcal{K}$ consists of:
\begin{enumerate}
    \item A carrier type $\text{Carrier}$.
    \item A view mapping $\text{view\_of} : \text{ISKSubtype} \to \text{Carrier}$.
    \item An observational equivalence relation $\approx$ on $\text{Carrier}$.
    \item A soundness property: operational equivalence in the substrate implies view equivalence:
    $$\forall t, u, \quad t \sim_{op} u \implies \text{view\_of } t \approx \text{view\_of } u$$
    \item A decoding/reconstruction mapping $\text{decode} : \text{Carrier} \to \text{ISKSubtype}$.
    \item Coherence axioms:
    \begin{itemize}
        \item $\text{decode}(\text{view\_of } t) \sim_{op} t$
        \item $\text{view\_of}(\text{decode } c) \approx c$
        \item $c_1 \approx c_2 \implies \text{decode } c_1 \sim_{op} \text{decode } c_2$
    \end{itemize}
\end{enumerate}
\end{definition}

\begin{definition}[Canonical ISAR Kernel]
\lean{ISAR.ISAR_Kernel}
\uses{ISAR.Kernel}
The identity kernel mapping the substrate directly to itself is the canonical presentation.
\end{definition}

\begin{definition}[Computable ISAR Kernel]
\lean{ISAR.ComputableISAR_Kernel}
\uses{ISAR.Kernel}
The computable kernel parametrized by normalization fuel.
\end{definition}

\begin{definition}[Optimal Computable Kernel]
\lean{ISAR.ComputableISAR_Kernel_Optimal}
\uses{ISAR.ComputableISAR_Kernel}
The computable kernel instantiated with optimal fuel computed from the term size for linearly-typed terms.
\end{definition}

\begin{definition}[Kernel Morphism]
\lean{ISAR.KernelHom}
\uses{ISAR.Kernel}
A morphism $f : K_1 \to K_2$ in the category $\mathcal{K}$ is a carrier mapping $\text{hom} : K_1.\text{Carrier} \to K_2.\text{Carrier}$ that:
\begin{itemize}
    \item Preserves the view representations: $\forall t \in \text{ISKSubtype}, \quad K_2.\text{view\_of } t \approx_2 \text{hom}(K_1.\text{view\_of } t)$
    \item Respects carrier equivalence: $\forall c_1, c_2, \quad c_1 \approx_1 c_2 \implies \text{hom}(c_1) \approx_2 \text{hom}(c_2)$
\end{itemize}
\end{definition}

\begin{definition}[Canonical Homomorphism]
\lean{ISAR.canonical_hom}
\uses{ISAR.KernelHom, ISAR.ISAR_Kernel}
For any kernel $K$, the canonical homomorphism into $\text{ISAR\_Kernel}$ is defined by the decoding mapping.
\end{definition}

\section{Universal Factorization and Terminality}

The uniqueness and terminality theorems are formalized and proved in \url{src/ISAR/KernelCategory.lean}. We prove that the canonical representation-free kernel is terminal, meaning that any admissible view can be uniquely decoded into it.

\begin{theorem}[Morphism Uniqueness]
\lean{ISAR.morphism_uniqueness}
\uses{ISAR.KernelHom, ISAR.canonical_hom}
Any structure-preserving morphism $f : K \to \text{ISAR\_Kernel}$ is observationally equivalent to the canonical decoding morphism:
$$\forall c \in K.\text{Carrier}, \quad f(c) \sim_{op} \text{decode } c$$
\end{theorem}

\begin{theorem}[ISAR Kernel Terminality]
\lean{ISAR.ISAR_Kernel_terminal}
\uses{ISAR.morphism_uniqueness, ISAR.ISAR_Kernel}
The quotient of the morphism space $\text{KernelHom } K\ \text{ISAR\_Kernel}$ modulo observational equivalence is a singleton:
$$\exists! [f] \in \text{KernelHom } K\ \text{ISAR\_Kernel} / \sim_{op}$$
Thus, $\text{ISAR\_Kernel}$ is a terminal object in the category $\mathcal{K}$ of semantic kernels.
\end{theorem}

This category-theoretic result is not a naming convention: it is an algebraic theorem stating that the Invariant Layer is the unique (up to isomorphism) cofinal quotient that retains only the information visible to observational equivalence. Any other view is mathematically guaranteed to factor uniquely through it, separating the dialect-specific representation from the coordinate-free invariant semantics.

\section{The Dialect Subsystem}

The Dialect structures in this section are verified in \url{src/ISAR/DialectKernel.lean}. A Dialect is an alternative abstract view over the substrate.

\begin{definition}[Dialect]
\lean{ISAR.Dialect}
A Dialect specifies:
\begin{enumerate}
    \item A type of dialect objects.
    \item A type of observations and observational equivalence.
    \item An evaluation mapping, an encoder to substrate terms, and a decoder from the substrate.
    \item A coherence law: evaluation commutes with encoding and decoding.
\end{enumerate}
\end{definition}
